\section{光的偏振}
光的干涉和衍射线性揭示了光的波动特性，而我们知道，光作为一种电磁波是一种横波，区分横波和纵波的关键差别在于振动方向是垂直还是平行传播方向，而在三维中，横波的振动方向将有一个平面的可能范围。具体的说，同样是这么一束光，它的光矢量究竟在垂直光线的平面的哪个方向上振动，上下？左右？或两者兼有之？这方面的特性我们称为\uwave{光的偏振}，而很显然的是，偏振必然只是横波独有的特性，因此在一定意义上，光的偏振现象确立了光的横波性。

\subsection{光的偏振态}
在光的传播过程中，如\xref{fig:完全偏振光}所示，假设光沿$z$轴传播，如果光矢量$\vb*{E}$仅在某个特定方向上振动，对于\xref{fig:水平偏振}这是水平的$x$轴方向，对于\xref{fig:垂直偏振}折射垂直的$y$轴方向，那么，我们就将这样的光称为\uwave{完全偏振光}，同时，我们将光矢量的振动平面，在这里分别是$zx$和$zy$平面，称为\uwave{偏振平面}。另外如果迎着光线看（左侧圆形区域的视角），完全偏振光的光矢量的振动始终保持在一条直线上，因此，完全偏振光有时也被称作\uwave{线偏振光}。为了在光线传播方向上也能表示出光的偏振方向，我们还会在光线上附加“圆点”或“竖线”以表示光矢量的振动方向，前者表示垂直纸面沿$x$轴的水平偏振，后者则表示纸面上垂直光线沿$y$轴的垂直偏振。

\begin{Figure}[完全偏振光]
    \begin{FigureSub}[水平偏振]
        \includegraphics[]{build/Chapter03G_02.fig.pdf}
    \end{FigureSub}\vspace{0.5cm}
    \begin{FigureSub}[垂直偏振]
        \includegraphics[]{build/Chapter03G_01.fig.pdf}
    \end{FigureSub}
\end{Figure}

就线偏振光这种迎着光线方向观察光矢量划出的轨迹的分类方式，假如光矢量随着光线传播的过程中，绕端点不断的匀速左旋或右旋，就会勾勒出圆形轨迹，这样的光称为\uwave{圆偏振光}，而较一般的说，假如光矢量勾勒出的是椭圆轨迹，这时就得到了\uwave{椭圆偏振光}。椭圆偏振光是最一般的情况，其“长轴和短轴相等”和“短轴为零”分别就给出了圆偏振光和线偏振光，不过反过来说，椭圆偏振光和圆偏振光也可以视为两个频率相同，振动方向垂直的线偏振光以不同或相同比例叠加的结果。总之，尽管三者的偏振方式略有不同，但它们都具有规律的偏振特性。

\uwave{自然光}，或者更普遍的说，以自发辐射为主的光源发出的光，是否是某种偏振光？答案是否定的，自然光并不是偏振光，从宏观上来看，自然光的光矢量的取向是随机的，在统计上有均等的概率向各个方向振动。这并不是说，自然光是一种圆偏振光，圆偏振光的光矢量虽然也在不断变化，但这种变化是具有确定规律的，自然光的偏振方向则是迅速且无规则的变化着的，这可以从普通光源的发光性质解释，在\xref{subsec:光源的发光机制}中我们提到过，普通光源的发出的光波实际是一段一段的，每一段光波来自不同的电子跃迁，因而，每一段光波虽然可能各自都具有规律的偏振状态，但彼此之间的偏振状态各不相同，从宏观上看就表现为光矢量无规律的变化了。

自然光的光矢量具有向各个方向振动的可能，但由于每一个光振动矢量都可以在两个互相垂直的方向上分解，因此，自然光的光矢量可以用两个，振动方向垂直，振幅相同，\empx{但没有恒定相位关系}的独立振动代替，由于两个独立振动的相位关系不断变化，它们的振动合成结果也会不断随机变化。基于这样的想法，\empx{自然光就可以视作两个振动方向垂直且振幅相等，但没有恒定相位关系的线偏振光的混合}。因此，自然光在偏振图示中常被表示为\xref{fig:自然光}的形态。
\begin{Figure}[自然光]
    \includegraphics{build/Chapter03G_03.fig.pdf}
\end{Figure}

自然光有时经过一些光学过程后，虽然仍然没有规律的偏振方向，但统计上观察到其光矢量在各个振动方向上出现的概率不再均等，某个方向上的振动多些，某个方向上的振动少些，我们将这种光称为\uwave{部分偏振光}，它可以视为完全偏振光和自然光的结合，图示如\xref{fig:部分偏振光}所示。
\begin{Figure}[部分偏振光]
    \begin{FigureSub}[水平部分偏振]
        \includegraphics[]{build/Chapter03G_05.fig.pdf}
    \end{FigureSub}\vspace{0.5cm}
    \begin{FigureSub}[垂直部分偏振]
        \includegraphics[]{build/Chapter03G_04.fig.pdf}
    \end{FigureSub}
\end{Figure}

\subsection{光偏振片}
现在的问题是，偏振光和自然光对我们人眼感知而言，有什么差异吗？答案是否定的，就像是自然光和人造白光源对在我们的感知而言都是“白光”，但是，自然光是连续光谱，人造白光源往往是特定几个波长的单色光的混合，由此可见，\empx{人眼无法感知光的全部的波动特性}，而偏振性也是我们无法感知的。但是，在某些天然或人造的材料存在一个特定方向，使得
\begin{itemize}
    \item 当光的偏振方向与之相垂直时，光会被强烈吸收。
    \item 当光的偏振方向与之相平行时，光会完全的通过。
\end{itemize}
这就好比水中放置的一个栅栏，如果鱼上下摆动就可以通过，如果鱼左右摆动就不能通过了。

这种对光的偏振具有选择性吸收的性质称为材料的\uwave{二向色性}，由这种材料制备的光学元件则称为\uwave{偏振片}，而偏振片能透过的光振动的方向称为的\uwave{偏振化方向}，这通过会标注在偏振片上。

偏振片的第一种用法是作为\uwave{起偏器}，\uwave{起偏}是指从自然光获得偏振光的过程，如\xref{fig:起偏器}，设通过偏振片前后的光强分别为$I_0$和$I$，由于自然光可以视为垂直和平行偏振化方向的两个完全偏振光的叠加，两者各占有自然光$I_0/2$的光强，而通过偏振片后，就只留下了平行偏振化方向的那一完全偏振光，因此通过偏振片后产生的偏振光的光强$I=I_0/2$就只有原先的一半了。
\begin{Figure}[起偏器]
    \includegraphics[]{build/Chapter03G_06.fig.pdf}
\end{Figure}

偏振片的第二种用法是作为\uwave{检偏器}，这研究的是偏振光再通过偏振片的结果，如\xref{fig:检偏器}所示
\begin{itemize}
    \item 旋转偏振片，使得其偏振化方向与偏振光平行，此时透射光最强。
    \item 旋转偏振片，使得其偏振化方向与偏振光垂直，此时透射光为零，称为\uwave{消光现象}。
\end{itemize}
\begin{Figure}[检偏器]
    \begin{FigureSub}[最亮]
        \includegraphics[]{build/Chapter03G_07.fig.pdf}
    \end{FigureSub}\vspace{0.5cm}
    \begin{FigureSub}[最暗]
        \includegraphics[]{build/Chapter03G_08.fig.pdf}
    \end{FigureSub}
\end{Figure}
通过检偏器，可以分辨偏振光的偏振方向，只要观察偏振片转到哪一方向透射光最亮即可。

\subsection{马吕斯定律}
现在的问题是，偏振化方向与偏振光平行则透射光最强，偏振化方向与偏振光垂直则透射光最弱。那么如果两者夹角$\alpha$介于平行$\alpha=0$和垂直$\alpha=\pi/2$之间，此时将会产生什么结果？

\begin{BoxLaw}[马吕斯定律]
    若偏振片的偏振化方向与偏振光的夹角为$\alpha$，则
    \begin{Equation}
        I=I_0\cos^2\alpha
    \end{Equation}
    其中$I_0,I$是通过偏振片前后的光强。
\end{BoxLaw}
\begin{Proof}
    设通过偏振片前后光矢量的振幅分别为$A_0,A$，而光强分别为$I_0,I$，我们将通过偏振片前的光矢量正交分解至平行偏振方向$A_0\cos\alpha$和垂直偏振方向$A_0\sin\alpha$，只有前者得到保留，故
    \begin{Equation}*
        A=A_0\cos\alpha
    \end{Equation}
    而光强和光矢量的振幅的关系是$I=A^2$和$I_0=A_0^2$，故
    \begin{Equation}*
        I=I_0\cos^2\alpha\qedhere
    \end{Equation}
\end{Proof}
马吕斯定律与上一小节是兼容的，当$\alpha=0$时$I=I_0$最亮，当$\alpha=\pi/2$时$I=0$最暗。

\subsection{反射和折射中的偏振现象}
马吕斯在1808年偶然从窗玻璃反射的太阳光中发现了光的偏振现象，进一步的研究表明，偏振光除了可以通过前面几小节的偏振片产生，当自然光在两种介质的界面上发生反射和折射时，反射光和折射光都会出现一定程度的偏振现象，即，反射光和折射光均为部分偏振光。
\begin{Figure}[反射和折射中的偏振现象]
    \includegraphics[]{build/Chapter03G_09.fig.pdf}
\end{Figure}
反射光和折射光的部分偏振的方向是不同的，如\xref{fig:反射和折射中的偏振现象}所示
\begin{itemize}
    \item 反射光的部分偏振，主要以垂直纸面的方向为主。
    \item 折射光的部分偏振，主要以平行纸面的方向为主。
\end{itemize}
这里所谓的“纸面”实质上，是入射光所在那一垂直界面的平面，有时也称为“入射平面”。

改变入射角，反射光和折射光的偏振化程度会相应发生变化，\empx{折射光在该过程中始终是部分偏振}，然而入射角在某些特定角度时，\empx{反射光则会变为完全偏振}，这由\uwave{布儒斯特定律}给出。
\begin{BoxLaw}[布儒斯特定律]
    当反射光和折射光垂直时，反射光将变为完全偏振光。
\end{BoxLaw}
定义反射光变为完全偏振光时的那一特殊入射角，定义为\uwave{布儒斯特角}，布儒斯特定律只告诉我们布儒斯特角对应的是反射光和折射光垂直的情形，并没有给出布儒斯特角的计算公式。

\begin{Figure}[布儒斯特角]
    \begin{FigureSub}[光由空气射向玻璃]
        \includegraphics[scale=0.85]{build/Chapter03G_10.fig.pdf}
    \end{FigureSub}\hspace{0.2cm}
    \begin{FigureSub}[光由玻璃射向空气]
        \includegraphics[scale=0.85]{build/Chapter03G_11.fig.pdf}
    \end{FigureSub}
\end{Figure}\vspace{0.2cm}

\begin{BoxFormula}[布儒斯特角]
    当入射角$i=i_B$时，反射光将变为完全偏振光
    \begin{Equation}
        \tan i_B=\frac{n_2}{n_1}
    \end{Equation}
    其中$i_B$即为布儒斯特角，光由介质$n_1$射向介质$n_2$。
\end{BoxFormula}
\begin{Proof}
    根据\fancyref{law:布儒斯特定律}，反射光变为完全偏振光时，反射光和折射光垂直，即
    \begin{Equation}&[1]
        i_B+r=\frac{\pi}{2}
    \end{Equation}
    根据\fancyref{law:光的折射定律}
    \begin{Equation}&[2]
        n_1\sin i_B=n_2\sin r
    \end{Equation}
    在\xrefpeq{2}代入\xrefpeq{1}
    \begin{Equation}&[3]
        n_1\sin i_B=n_2\sin\qty(\frac{\pi}{2}-i_B)=n_2\cos i_B
    \end{Equation}
    这就得到
    \begin{Equation}*
        \tan i_B=\frac{n_2}{n_1}\qedhere
    \end{Equation}
\end{Proof}
正是因为反射光是具有一定的偏振性的，因此在摄影中有时会在镜头前附加偏振镜，当偏振镜被旋转到某些特定角度时，就可以将水面的反光消除，使得水面下的景象清晰的呈现出来。